\documentclass[11pt,a4paper]{article}
\usepackage[margin=1in]{geometry}
\usepackage{amsmath,enumitem,xr}
\usepackage[hidelinks]{hyperref}
\externaldocument{main}
\title{Response to the Supplied Editorial Concerns}
\author{Karem Abdelgalil}\date{6 October 2026}
\begin{document}\maketitle
This revision preserves the manuscript's dictionary, projection, alignment, completion, and interaction structure while correcting claims that exceeded the specified models. This letter addresses the three scientific concerns supplied with the manuscript and the request to consolidate scope statements. It does not imply that every other editorial requirement has been verified.

\section*{1. The central result was a reparameterization without a valid observation model}
We agree that adding a preferred phase to a transfer equation does not derive it from a receptive gate. Section~\ref{sec:phase} now separates transfer phase $\gamma$ from receptive center $\alpha$. A positive real gate changes magnitude, while the distinct factor $e^{-i\gamma}$ rotates the carrier. The observation mean is $a g(\psi-\alpha-\Omega\tau)e^{-i(\gamma+\Omega\tau)}$, with explicit independent real and imaginary Gaussian errors of variance $s^2$. The likelihood and information use this consistent variance convention.

Proposition~\ref{prop:orbit} gives the full three-parameter equivalence and its null direction. The phase-only rank-one result concerns $(\gamma,\tau)$ and does not identify the receptive center. Raw simultaneous phase contrast is a probe variable, rather than an equation containing these physiological parameters. A two-band observation with shared transfer phase plus a gain scan identifies the three parameters locally on a calibrated phase branch. The result is classical identifiability analysis of an explicit observation, not a new universal mathematical theorem. Common drive is modeled by sums of analytic signals, not by adding complex signals to real angles.

\section*{2. The paper combined a short identification result with a definitional report framework}
The author has elected to retain the framework, while explicitly distinguishing its postulates from derived consequences. Section~\ref{sec:limits} defines report-status vocabulary once and states that the model has no experience variable. The image test is corrected to scalar-readout nonorthogonality, with a separate item-specific amplitude condition. The gate is positive even at opposite phase. Proposition~\ref{prop:delivery} derives the exact threshold window, and Theorem~\ref{thm:inter} establishes a nonzero sender-dependent receiver contribution and an exact source-content contrast under the isolated-path model. Corollary~\ref{cor:completion} states threshold feasibility and operational first passage without claiming that supported content must finish.

The communication theorem follows from a specified transfer model, and the threshold corollary follows from an operational criterion; neither is promoted as empirical proof of brain communication. Carrier compensation does not preserve arbitrary transient onset times. The sensitivity corollary now follows only from an explicit Gaussian readout with downstream noise. Projection geometry and monotonicity alone do not force a positive global regression interaction. The bilinear contrast detects cross-path dependence against separable alternatives but cannot identify its neural location. The slower trace is retained as a learning postulate whose simulated split is imposed.

Section~\ref{sec:reciprocal} now gives both levels a common reciprocal input--output model. Its time-domain echo series and stationary carrier response contain the same directed map products. Theorem~\ref{thm:reciprocal} derives item-specific return support; Corollary~\ref{cor:recurrence} derives conduction plus round-trip waiting for an explicitly positive invariant channel. Reciprocity alone preserves two single-band parameter-null directions, while independently timed onsets and resolved returns can supply transverse timing information. These are conditional consequences with full proofs, not a claim that the elementary geometry has discovered a physiological latency. One-way transfer remains a limiting case. Whether the journal considers the resulting synthesis sufficiently novel remains an editorial judgment.

\section*{3. Simulation reproducibility, unexplained coverage, and missing code}
The source package includes executable models, independent verification, fixed seeds, requirements, figures, per-realization results, a one-command numerical reproduction driver, and Supplementary Methods. The original supplied simulations are preserved and rerun separately; their source-mixing checks are explicitly distinguished from the revised recovery experiment.

The explanation that regressor correlation itself invalidates the OLS standard error is removed: OLS covariance already accounts for that correlation. The historical 52/60 coverage result is retained as provenance without an unsupported diagnosis. A revised 2,000-realization exact-feature experiment covers in 1,900/2,000 Student-$t$ intervals. With feature noise, naive coverage falls to 47/2,000. A fully implemented 2,000-resample pairs bootstrap covers in 57/60 exact-feature realizations; analytic intervals cover the same number on that subset. Its uncertainty is reported and superiority is not asserted.

The temporal rival comparison now states its actual feature sets. A quadratic including the imposed delayed product performs almost as well as the delayed bilinear predictor. This limits mechanistic interpretation. The supplement specifies numerical settings and distinguishing implementation checks from neural validation.

The author supplies the public project repository \url{https://github.com/lolzguycoder/esj-phase-connectivity} and archive DOI \href{https://doi.org/10.5281/zenodo.23165931}{10.5281/zenodo.23165931}. They are included in the final code-availability statement. Current archive contents and their correspondence to this revision could not be independently verified, so exact-release deposition is not asserted.

\section*{Expository revision following the supplied external review}
The complete proofs and eight-section structure are retained. A conceptual roadmap distinguishes the measurement and content branches; a reading guide identifies independent sections. The abstract is shortened and organized around problem, approach, results, and limits. Revised code and figures use $\alpha$ consistently for receptive center. The signal-detection and bilinear models are introduced as complementary behavioral and geometric extensions. A worked two-coordinate learning example and an unexcited-direction corollary make the standard mean-recursion result concrete, with the standard stochastic-approximation literature cited.

Table~\ref{tab:reject} is rebuilt with ragged-right columns and explicit row spacing. The source already used distinct upstream and downstream noise symbols; they are now explicitly labeled in the display. The wPLI range, two-band design, canceling multi-band derivative example, linear source contrast, and basis invariance are stated. The comparison with Dowdall and Vinck is based on their published bidirectional interference result and does not imply that the present identifiability model derives that mechanism. The Nolte reference already spells M\"uller consistently.

\section*{Proof expansion and readability}
All fifteen main-text proofs are expanded to give their constructions, intermediate equalities, and boundary cases. The interaction-rank proof is also expanded in the supplement. The completion corollary explicitly assumes a continuous gain trajectory to ensure that a finite crossing-time infimum is attained. Learning assumptions now state the initial-error integrability needed by the expectation argument. Explanatory transitions distinguish structural support, item delivery, threshold crossing, interaction rank, and expected learning. The original result scope is retained.

\section*{Additional scientific development}
The eight-section structure is retained. The revised projection result is extended by uncertainty bounds that distinguish certified delivery from uncertain map estimates. The effective bilinear rank criterion is promoted into the main paper to specify which receiver directions a two-path interaction can add. A gain-weighted excitation proposition derives expected translation-error dynamics and identifies errors that restricted training cannot correct. An integrated experimental opportunity links these conditional results to held-out content contrasts, timing calibration, complement-rank measurement, and restricted-subspace training. Existing simulations remain synthetic; new randomized algebra checks verify the additional results. No established biological breakthrough is claimed.

\section*{Latest eight-point editorial revision}
The retarded operators are closed into a common convolutional feedback model; the sharp discrete-delay loop is its instantaneous-processing limit, with the point-mass and continuity qualifications stated. A new general-return proposition gives exact matrix first passage, tail certificates, transient-overshoot examples, and a finite support test beyond the invariant-channel corollary. A Gaussian innovation-readout corollary links the interaction-rank result to discriminability while requiring class-dependent content and fixed calibrated noise. Thus added rank is an opportunity, not an automatic behavioral effect.

The former Sections 4.1 and 4.2 are merged into a single latency discussion. Learning, interaction, gain, and noise symbols are defined locally before their equations. Dehaene's name and amplitude notation are checked; maximum projected amplitude $A_v$ remains distinct from instantaneous $a_v$, and no bare ``av'' or ``1This'' occurs in the revised source. A 300-realization paired comparison supplies pointwise MSE-difference intervals and four-test Holm adjustments, without treating temporal samples as independent replicates. The title and Section 4 heading now emphasize phase-gated interregional transfer rather than an undeveloped unreported-processing theory.

Public-data feasibility is added to the manuscript and Supplement S16, with a separate source-linked data review. CCEP metadata are inspected for timing fields, electrode labels and bilateral coverage; Allen datasets are assessed for population features, stimulus labels and behavior. No neural waveform is fitted and no inspected resource is asserted to validate the complete causal model. Publication remains a theoretical and synthetic contribution with a concrete empirical route.

\section*{4. Scope statements and submission status}
The proposed prefrontal/contralateral application names specific calibrated populations rather than treating a hemisphere as a single region or assigning conscious status by anatomy. Published no-report prefrontal decoding and callosal-delay studies motivate the application but do not validate the proposed loop. A new reproduction subsection specifies all reciprocal maps, gains, units, delays, threshold cases, and independent random recurrence checks. The full driver regenerates the reciprocal figure and numerical output alongside all prior analyses.
Broad consciousness claims and repeated assurances have been replaced by narrower statements attached to the affected equations. The abstract reports the actual conditional model results and synthetic evidence. No biological validation or breakthrough is claimed. The author must verify the revised source, check archive-to-revision correspondence, confirm reference metadata and journal-specific requirements before submission.
\section*{Final detailed review: corrections and extensions}
The reciprocal theorem and proof consistently use $L_{PQ}$, $L_{QP}$ and $A=L_{PQ}L_{QP}$; the item Krylov space uses the same map. The Jacobian and offset information matrix were already present in source and have been checked in the rendered revision. The notation table now separates $U,V$, the bilinear map, memory kernel and interaction coefficient. Existing correct author spellings and accented names are retained. The wPLI is explicitly identified as unsigned.

The non-normal suggestion required a mathematical correction: an induced norm below one excludes vector-norm transient amplification. A new fully proved proposition instead extends sharp-loop completion certification to the broader Schur-stable case using a Lyapunov metric, with an independently checked non-normal numerical example. The proposed upper-triangular example with source $e_1$ and readout $e_2$ would have zero readout at every return; the revised example uses source $e_2$ and readout $e_1$.

The revision states readout/noise and alternative-learning limits, clarifies across-condition span versus single-trial attainability, explains independent echo timing, adds alias calibration and expands the harmonic/filter interpretation. Precision planning supplies formulas and explicitly hypothetical examples rather than invented biological sample minima. A minimal experiment, formal comparison table, Granger discussion, symbol glossary, gate-calibration note and innovation-geometry figure are included. Captions report simulation truths, initial errors and the paired-comparison seed. A shorter abstract and sharper positioning retain phase-gated transfer as the main scope, with report as motivation. The paper remains one integrated study, as requested.

The expanded public-data review includes Dryad frontal subspace recordings, DANDI releases for hippocampal and cortical recordings and for cortical stimulation, IBL decision-task data, and a human PFC--M1 study with restricted additional-data access. Public DANDI version metadata were retrieved; no neural waveforms were fitted. These resources strengthen component-test feasibility and related-work positioning. None documents the complete joint intervention and calibration design, so the revision does not present them as empirical validation or proof of a breakthrough.
\section*{Additional nervous-system signaling request}
The revision keeps one paper and the Enosh class. Theorem~\ref{thm:network} generalizes the retarded operator to multiple directed routes, deriving item-specific supported arrival and cancellation-aware threshold certificates with a complete proof. Proposition~\ref{prop:dispersion} bounds processing dispersion and exhibits an exact conduction--processing ambiguity even with broadband observations. The new verifier independently compares network implementations and tests signed-route and dispersion bounds. Supplement S19 records its seed, settings, and exclusions. A sensorimotor public-data route is added after the existing experimental plan, with raw-data validation explicitly deferred. Projection establishes a modeled source-dependent effect under calibrated conditions; it does not identify a missing physiological delay or a conscious/subconscious anatomical division.

\section*{Final network review: minor revisions}
The paper remains one study in the Enosh template. A fully proved weighted-network lemma relaxes the unweighted row condition and recovers the 2 versus 0.4 two-edge example. The comparison spectral-radius condition remains sufficient for signed networks. The dispersion discussion now establishes asymptotic sharpness, gives non-uniform kernels, and specifies independent constraints that break conduction--processing equivalence. Learning updates are explicitly connected to frozen network edge maps, with stability rechecked after training.

Section 7.3 now gives exact assumed-effect paired-$t$ power calculations, an equivalence-power calculation for a near-null contrast, and calibrated rejection criteria. The new official-manifest and published-session-table audit distinguishes 49 files from 16 simultaneous-area candidate sessions across the four principal areas, including only nine PMd--dlPFC sessions and no direct dlPFC--M1 pair. Those sessions come from two animals, so the release cannot be presented as adequately powered validation of the full framework. Raw download access errors and the absence of a new biological fit are stated. The chi symbol is spelled out in the notation table; Supplementary Figure S9 is included. S20 and executable calculations document the added checks.
\end{document}
